% ============================================================
%  part4/ch19-sobolev.tex  —  Chapter 19: Sobolev Spaces
% ============================================================

\chapter{Sobolev Spaces and Function-Space Regularity}
\label{ch:sobolev}

\epigraph{%
  The geometry does not determine the analysis.\\
  The analysis determines which geometry is legitimate.%
}{---}

\begin{WhyMatters}
  This chapter establishes the methodological chain that
  governs everything from here forward:
  \[
    \text{Field Equations}
    \Rightarrow \text{Regularity}
    \Rightarrow \text{Sobolev Category}
    \Rightarrow \text{Admissibility Sheaf}
    \Rightarrow \text{Cohomology.}
  \]
  The admissibility sheaf is not chosen because sheaves are
  elegant.
  It is derived because the PDE theory demands a specific
  function space.
\end{WhyMatters}

\section{Sobolev spaces}

\begin{definition}[Sobolev space $H^k$]
  \label{def:sobolev}
  Let $\Omega \subset \mathbb{R}^n$.
  The \emph{Sobolev space} $H^k(\Omega)$ is
  \[
    H^k(\Omega)
    = \{u \in L^2(\Omega) : D^\alpha u \in L^2(\Omega),\;
      |\alpha| \leq k\},
  \]
  with norm
  \[
    \|u\|_{H^k}^2 = \sum_{|\alpha|\leq k} \|D^\alpha u\|_{L^2}^2.
  \]
\end{definition}

\begin{definition}[RSVP Sobolev state space]
  \label{def:rsvp-sobolev}
  For an open set $U \subseteq X$, define
  \[
    \mathcal{H}^k_{\mathrm{RSVP}}(U)
    = H^k(U) \times H^k(U; TU) \times H^k(U) \times H^k(U).
  \]
  An element is $s_U = (\Phi_U, \mathbf{v}_U, S_U, \ell_U)$.
  The physical lamphron density is $L_U = \ell_U^2$.
\end{definition}

\begin{proposition}[Positivity of lamphron density]
  \label{prop:lamphron-positive}
  \statusproven{}
  For every $\ell_U \in H^k(U)$, the physical lamphron field
  $L_U = \ell_U^2$ satisfies $L_U(x) \geq 0$ a.e.
\end{proposition}

\begin{proof}
  For any real $a \in \mathbb{R}$, $a^2 \geq 0$.
  Since $\ell_U$ is real-valued a.e., $L_U(x) = \ell_U(x)^2 \geq 0$ a.e.
\end{proof}

\section{Regularity requirements from the field equations}

\begin{proposition}[Allen--Cahn regularity]
  \label{prop:ac-regularity}
  \statusproven{}
  For the Allen--Cahn equation $\partial_t\Phi = \Delta\Phi + f(\Phi)$:
  weak solutions require $\Phi \in H^1$;
  classical solutions require $\Phi \in H^2$.
\end{proposition}

\begin{proof}
  The weak formulation pairs $\partial_t\Phi = \Delta\Phi + f(\Phi)$
  against test functions $\psi$:
  $\int_U \Delta\Phi\,\psi\,dx = -\int_U \nabla\Phi\cdot\nabla\psi\,dx$.
  This requires $\nabla\Phi \in L^2$, i.e., $\Phi \in H^1$.
  For $\Delta\Phi \in L^2$, one needs $\Phi \in H^2$.
\end{proof}

\begin{proposition}[Cahn--Hilliard regularity]
  \label{prop:ch-regularity}
  \statusproven{}
  For the Cahn--Hilliard equation $\partial_t L = -\Delta^2 L + \Delta f(L)$:
  weak energy solutions require $L \in H^2$;
  classical solutions require $L \in H^4$.
\end{proposition}

\begin{proof}
  The free energy $\int|\nabla L|^2 dx$ controls $\|\nabla L\|_{L^2}$,
  requiring $L \in H^1$.
  The chemical potential $\mu = -\Delta L + f(L)$ requires $\Delta L \in L^2$,
  i.e., $L \in H^2$ for weak solutions.
  Classical pointwise interpretation of $\Delta^2 L$ requires $L \in H^4$.
\end{proof}

\begin{StatusBox}{Proven Framework}
  The Sobolev index $k$ for the full coupled RSVP system is
  \statusopen{}.
  It is determined by the highest-order term in the equations,
  which depends on whether the lamphron evolution is
  Allen--Cahn-type ($k=1$ sufficient) or Cahn--Hilliard-type
  ($k=2$ minimum).
  The well-posedness chapter settles this.
\end{StatusBox}

\begin{MathEcho}
  This chapter introduced:
  $H^k(\Omega)$, $\mathcal{H}^k_{\mathrm{RSVP}}(U)$,
  $L_U = \ell_U^2 \geq 0$.
  The Sobolev index $k$ will determine the target category
  of the admissibility sheaf in Chapter~\ref{ch:sheaf}.
\end{MathEcho}
